%!TEX root = ../hutbthesis_main.tex

\chapter{总结与展望}
\section{总结}
本研究主要围绕肌肉驱动的生物力学人机交互系统展开，重点做了系统的设计和实现工作，我们以方向盘控制作为核心应用场景，一步步完成了从模型搭建、耦合驱动实现，再到系统功能测试的全部流程，没有遗漏任何关键步骤。

本系统借助 User-in-the-Box 开源框架完成搭建工作，框架可以为研究提供可用的人体上肢肌肉骨骼模型，模型内部包含胸廓、脊柱、肩部、肘部、腕部等多个关节结构，关键肌肉发力组件与关节力矩计算端口也被集成在模型中，肌肉的激活信号可以被转化为关节转动需要的动力，这一过程为肌肉驱动操作提供了基础条件，也让后续的控制流程可以稳定开展。

研究人员通过仿真文件将上肢模型与方向盘模型进行固定连接，肌肉激活信号、关节力矩、方向盘转动状态被连接成一条完整的控制链路，上肢模型做出转动动作时，方向盘会跟随上肢同步运动，两者可以保持一致的运动状态。

模型搭建完成后需要进行合理性检验，工作人员从结构完整程度、动力响应效果、运动限制情况等角度开展测试，测试结果表明，模型结构完整稳定，运动过程不会出现关节超出限制、姿态异常、抖动偏移等问题，可以满足仿真使用的基本要求。

系统功能测试环节设置了三种不同任务，分别为遥控驾驶任务、控制精度任务、小物体触碰任务，这些任务可以检验连续操作、精准定位、狭小空间交互等不同场景下的系统表现。

测试得到的数据显示，模型运行过程平稳顺畅，控制效果稳定可靠，方向盘与上肢的联动没有明显延迟，指尖定位精度与小物体触碰成功率可以达到人机交互的使用标准，这一结果说明，肌肉与方向盘耦合驱动系统具备实际使用价值，能够在模拟驾驶、精细操作等多种场景中完成对应的交互任务。
\section{不足与展望}
我们开展的这项研究，目前还存在一些不足之处，这些问题还需要在之后的工作中进一步改进和完善。

首先在模型方面，目前我们使用的上肢模型，其肌肉驱动的逻辑还比较简单，没有考虑到肌肉疲劳、不同人身体状况不一样等因素，这些因素都会影响肌肉发力的效果，所以这个模型和真实人体肌肉的复杂运动情况还有一定的差距；另外，方向盘模型也只是实现了基础的固定绑定，没有加入真实方向盘所具备的阻尼、自动回正等物理特性，和我们实际开车时接触到的方向盘体验还有区别。

然后在控制方面，当前系统使用的是固定的按键控制方式，没有结合人体真实的肌电信号来进行闭环控制，所以人和机器之间的交互还不够自然，还有提升的空间；同时，自动评测模式也只能应对提前设置好的测试场景，不能应对复杂的工作情况和多目标的任务，扩展能力还比较弱。

这次的系统设计存在以上问题，后续将在以下方面进行改进。

第一，进一步完善生物力学模型，加入更细致的肌肉参数和人体生理方面的限制条件，让这个模型和真实人体的贴合度更高，更符合实际情况；

第二，优化方向盘的物理模型，增加阻尼、回正力矩等相关参数，让方向盘的使用体验更接近真实开车时的感觉，提升交互的真实性；

第三，添加肌电信号采集相关的设备，实现通过人体真实肌肉信号来进行闭环控制，让人和机器的交互更自然、更有沉浸感；

第四，增加更多的测试场景，比如复杂的路况、需要同时完成多个目标的交互场景等，不断完善系统的稳定性和适应能力，让肌肉驱动的生物力学人机交互技术，能够真正运用到驾驶辅助、康复训练等实际场景中去。
